\chapter{Glosario}
\label{chap:glosario}

Para la correcta comprensión del dominio del problema, se definen a continuación los conceptos económicos, financieros y técnicos fundamentales utilizados en el diseño del sistema:

% ==========================================
% SECCIÓN: TÉRMINOS FINANCIEROS Y DE TRADING
% ==========================================

\section*{Conceptos Financieros y de Trading}

\begin{description}
    \item[Exchange (Casa de Intercambio)] Plataforma digital que facilita la compraventa de activos financieros, como criptomonedas, actuando como intermediario y proveedor de liquidez y datos de mercado. En este proyecto, se integra con la API REST de Binance para la obtención de datos históricos y precios en tiempo real.

    \item[OHLCV] Acrónimo de \textit{Open, High, Low, Close, Volume}. Es el formato técnico estándar para representar la acción del precio de un activo financiero durante un intervalo de tiempo específico (una vela o \textit{candlestick}):
    \begin{itemize}
        \item \textbf{Open:} Precio al inicio del intervalo.
        \item \textbf{High:} Precio máximo alcanzado durante el intervalo.
        \item \textbf{Low:} Precio mínimo alcanzado durante el intervalo.
        \item \textbf{Close:} Precio al cierre del intervalo.
        \item \textbf{Volume:} Cantidad total del activo negociada durante el intervalo.
    \end{itemize}
    En el sistema, esta estructura se modela mediante la entidad JPA \texttt{Vela}.

    \item[Wallet (Billetera Digital)] Entidad encargada de custodiar y gestionar el capital del usuario. En el contexto de la aplicación, se implementa un \textit{modelo de doble saldo} para distinguir entre:
    \begin{itemize}
        \item \textbf{Balance Real:} Patrimonio total del usuario.
        \item \textbf{Balance Disponible:} Capital libre para nuevas operaciones (excluye fondos ya asignados a estrategias activas).
    \end{itemize}
    Esta separación previene sobregiros y facilita la gestión de riesgo.

    \item[Ledger (Libro Mayor)] Registro contable inmutable y secuencial donde se anotan todas las transacciones financieras. Su función es garantizar la integridad de los datos y permitir la auditoría de movimientos como:
    \begin{itemize}
        \item \textbf{RESERVED:} Fondos bloqueados para una estrategia.
        \item \textbf{COMMITTED:} Fondos en órdenes activas de trading.
        \item \textbf{REALIZED\_PnL:} Ganancia o pérdida realizada al cerrar una posición.
        \item \textbf{FEES:} Comisiones cobradas por el exchange.
    \end{itemize}
    En el sistema, cada movimiento genera una entrada de tipo \texttt{LedgerEntry}, permitiendo reconciliación en cualquier momento.

    \item[Estrategia de Trading] Conjunto de reglas lógicas y matemáticas que determinan de forma autónoma cuándo comprar o vender un activo. Estas reglas se implementan como subclases Python que heredan de \texttt{BaseStrategy}, permitiendo:
    \begin{itemize}
        \item Definición de indicadores personalizados.
        \item Condiciones de entrada (should\_buy).
        \item Condiciones de salida (should\_sell).
        \item Cálculo de stop-loss y take-profit.
    \end{itemize}

    \item[Indicador Técnico] Cálculo matemático derivado del precio y/o volumen histórico de un activo, utilizado para predecir movimientos futuros del mercado. Ejemplos clave implementados:
    \begin{itemize}
        \item \textbf{RSI (Relative Strength Index):} Oscilador que mide la velocidad y el cambio de movimientos de precios (rango 0-100). RSI < 30 sugiere sobreventa; RSI > 70 sugiere sobrecalentamiento.
        \item \textbf{SMA (Simple Moving Average):} Promedio aritmético de los últimos $n$ precios de cierre.
        \item \textbf{EMA (Exponential Moving Average):} Media móvil que da más peso a observaciones recientes.
        \item \textbf{MACD (Moving Average Convergence Divergence):} Indicador de impulso basado en diferencia de dos EMAs.
    \end{itemize}
    El sistema persiste estos valores en la entidad \texttt{IndicadorTecnico} para análisis posterior.

    \item[Paper Trading] Modalidad de simulación bursátil que permite operar con dinero ficticio en un entorno que refleja condiciones de mercado real (precios reales de Binance, pero sin dinero real comprometido). El sistema soporta este modo mediante el servicio \texttt{PaperTradingService}, permitiendo:
    \begin{itemize}
        \item Validar estrategias sin riesgo financiero.
        \item Educar usuarios en operativa antes de riesgo real.
        \item Recolectar datos de backtesting confiables.
    \end{itemize}

    \item[Backtesting] Proceso de validación histórica de una estrategia de trading. Consiste en simular la ejecución de una estrategia sobre datos de mercado pasados, sin estar conectado a un exchange real. Permite:
    \begin{itemize}
        \item Evaluar el rendimiento potencial antes de invertir capital real.
        \item Identificar estrategias débiles temprano.
        \item Optimizar parámetros de la estrategia.
    \end{itemize}
    En el sistema, el motor \texttt{engine\_backtest.py} itera sobre velas históricas simulando entry/exit según la lógica de la estrategia.

    \item[Lookahead Bias] Error crítico en backtesting donde el algoritmo ``ve el futuro''. Por ejemplo, usar el precio de cierre de mañana para decidir si comprar hoy. Esto infla artificialmente el rendimiento. El sistema lo evita mediante \textit{split temporal}, donde el modelo se entrena solo con datos históricos que existían en el momento de la predicción.

    \item[PnL (Profit and Loss)] Término que hace referencia a las ganancias o pérdidas netas obtenidas en una operación o en un período determinado. Se calcula como:
    $$\text{PnL} = \text{Precio\_Salida} - \text{Precio\_Entrada} - \text{Comisiones}$$
    En el sistema, se registra inmediatamente en el Ledger al cerrar una posición.

    \item[Win Rate] Porcentaje de operaciones rentables sobre el total de operaciones ejecutadas. Fórmula:
    $$\text{Win Rate} = \frac{\text{Número de Trades Ganadores}}{\text{Número Total de Trades}} \times 100\%$$
    Un win rate de 60\% significa que 6 de cada 10 trades son ganadores (no necesariamente lucrativo si pérdidas son muy grandes).

    \item[Max Drawdown] Pérdida cumulativa máxima experimentada desde un pico anterior de patrimonio. Mide el riesgo de una estrategia:
    $$\text{Max Drawdown} = \frac{\text{Capital\_Mínimo} - \text{Capital\_Pico}}{\text{Capital\_Pico}} \times 100\%$$
    Un max drawdown de -15\% significa que en el peor momento, el usuario vio su capital reducido 15\% desde su pico.

    \item[Profit Factor] Relación entre ganancias totales y pérdidas totales:
    $$\text{Profit Factor} = \frac{\text{Suma de PnL Positivos}}{\text{Suma Absoluta de PnL Negativos}}$$
    Un profit factor > 1.5 generalmente se considera aceptable.

\end{description}

% ==========================================
% SECCIÓN: TÉRMINOS TÉCNICOS DE MACHINE LEARNING
% ==========================================

\section*{Conceptos de Machine Learning}

\begin{description}
    \item[Feature Engineering] Proceso de transformación y combinación de datos brutos en variables (\textit{features}) que sean útiles para un modelo predictivo. En este proyecto, el \texttt{feature engineering dinámico} permite que cada estrategia defina sus propios features sin código duplicado, adaptando:
    \begin{itemize}
        \item Indicadores técnicos a calcular.
        \item Períodos de lookback.
        \item Normalizaciones o escalados especializados.
    \end{itemize}

    \item[Label (Etiqueta)] Variable de salida en problemas de clasificación supervisada. En este proyecto:
    \begin{itemize}
        \item \textbf{+1:} Señal de compra identificada por la estrategia.
        \item \textbf{-1:} Señal de venta identificada por la estrategia.
        \item \textbf{0:} Posición neutral (no hay señal).
    \end{itemize}
    Las labels se generan automáticamente mediante el método \texttt{get\_label(row, next\_row)} de la estrategia, permitiendo entrenamiento supervizado.

    \item[Desbalance de Clases] Situación donde las clases de un dataset de clasificación no están equitativamente distribuidas. En trading, es muy común (ej., 71\% clase neutra, 18\% compra, 11\% venta; véase Capítulo \ref{chap:desarrollo}, subsección \ref{subsec:desbalance_labels}). El desbalance severo causa que modelos naive predigan solo la clase mayoritaria. El sistema lo aborda con:
    \begin{itemize}
        \item Métricas ponderadas (\texttt{average='weighted'}).
        \item \texttt{class\_weight='balanced'} en XGBoost y Random Forest.
        \item Documentación explícita de distribución en metadata.
    \end{itemize}

    \item[Split Temporal] División train/test que respeta la causalidad temporal. A diferencia de split aleatorio, en series temporales es crítico entrenar solo con datos históricos y testear en datos más recientes:
    $$\text{Train:} t_1 \ldots t_{70\%} \quad \text{Test:} t_{70\%+1} \ldots t_{100\%}$$
    Esto evita lookahead bias y valida la generalización real.

    \item[Warmup Period] Número de observaciones iniciales descartadas antes de comenzar a usar indicadores técnicos o labels. Esto es necesario porque indicadores como RSI requieren al menos $n$ períodos históricos para poder calcularse. El sistema automáticamente:
    \begin{itemize}
        \item Identifica el máximo período de todos los indicadores (ej., 26 para MACD).
        \item Descarta las primeras 26 velas del DataFrame.
        \item Asegura que todas las features sean válidas (no NaN).
    \end{itemize}

    \item[Overfitting] Problema donde un modelo memoriza detalles específicos del dataset de entrenamiento en lugar de aprender patrones generalizables. En trading, es especialmente peligroso porque:
    \begin{itemize}
        \item Datos históricos pueden contener patrones espurios que no se repiten.
        \item Un modelo overfitted parece excelente en backtesting pero falla en operativa real.
    \end{itemize}
    El proyecto lo mitiga con split temporal, dropout en redes neurales, y validación en conjunto de test no visto.

    \item[Accuracy] Proporción de predicciones correctas sobre el total:
    $$\text{Accuracy} = \frac{\text{Predicciones Correctas}}{\text{Total de Predicciones}}$$
    Por sí sola, es engañosa en desbalance severo (un clasificador que siempre predice la clase mayoritaria puede tener 70\% accuracy).

    \item[F1 Score (Ponderado)] Media armónica entre Precision y Recall:
    $$F1 = 2 \cdot \frac{\text{Precision} \times \text{Recall}}{\text{Precision} + \text{Recall}}$$
    La versión ponderada (\texttt{average='weighted'}) considera la distribución de clases, siendo más realista en datasets desbalanceados.

    \item[Matriz de Confusión] Tabla que desglosa predicciones correctas e incorrectas por clase. Para un problema binario:
    \begin{itemize}
        \item \textbf{TP (True Positive):} Predicción positiva correcta.
        \item \textbf{FP (False Positive):} Predicción positiva incorrecta (falsa alarma).
        \item \textbf{TN (True Negative):} Predicción negativa correcta.
        \item \textbf{FN (False Negative):} Predicción negativa incorrecta (falta).
    \end{itemize}

    \item[Hyperparámetros] Parámetros de configuración del algoritmo de ML que no se aprenden durante el entrenamiento, sino que se especifican antes. Ejemplos:
    \begin{itemize}
        \item \textbf{XGBoost:} \texttt{n\_estimators=150, learning\_rate=0.05, max\_depth=6}.
        \item \textbf{Red Neuronal:} \texttt{epochs=50, batch\_size=64, learning\_rate=0.001, dropout\_rate=0.3}.
    \end{itemize}
    El sistema permite inyectar hiperparámetros desde Java hacia Python para experimentación sin cambiar código.

    \item[Dropout] Técnica de regularización en redes neuronales donde se desactivan aleatoriamente neuronas durante el entrenamiento, reduciendo dependencias co-adaptadas. Típicamente, \(\text{dropout\_rate} = 0.3\) (desactiva 30\% de neuronas en cada paso).

    \item[Metadata] Información sobre el entrenamiento de un modelo, como:
    \begin{itemize}
        \item Strategy utilizada, símbolo, timeframe.
        \item Rango temporal del dataset.
        \item Número de features y sus nombres.
        \item Métricas de validación (accuracy, F1, confusion\_matrix).
        \item Hyperparámetros utilizados.
    \end{itemize}
    Esto permite reproducir exactamente qué configuración produjo qué resultado.

\end{description}

% ==========================================
% SECCIÓN: TÉRMINOS TÉCNICOS DE ARQUITECTURA E IPC
% ==========================================

\section*{Conceptos de Arquitectura e Inter-Process Communication}

\begin{description}
    \item[IPC (Inter-Process Communication)] Mecanismo para intercambiar datos entre procesos distintos corriendo en el mismo o diferente máquina. En este proyecto:
    \begin{itemize}
        \item Proceso Java orquestador (PID principal).
        \item Procesos Python hijo (engine\_fetch.py, engine\_train.py, engine\_rt.py).
        \item Comunicación vía stdin/stdout con framing MessagePack.
    \end{itemize}

    \item[Protocolo IPC v1.0] Contrato explícito que define la estructura de todos los mensajes intercambiados:
    \begin{verbatim}
    {
      "protocol_version": "1.0",
      "message_type": "FETCH_REQUEST|TRAIN_RESPONSE|...",
      "correlation_id": "uuid-v4",
      "payload": { ... }
    }
    \end{verbatim}
    Permite evolución futura sin romper compatibilidad.

    \item[Framing Length-Prefixed] Técnica de serialización donde cada mensaje está precedido de su longitud en bytes (típicamente 4 bytes big-endian):
    $$\text{[4 bytes length]} + \text{[body]}$$
    Esto evita problemas de sincronización en tuberías: el receptor sabe exactamente cuántos bytes esperar.

    \item[MessagePack] Formato binario de serialización de datos más compacto que JSON (~30-50\% menor tamaño). Soporta tipos de datos complejos (arrays, maps) sin la verbosidad de JSON. En el proyecto se usa con framing length-prefixed para máxima eficiencia.

    \item[Fallback Legacy] Capacidad de un protocolo nuevo de aceptar mensajes en formato antiguo para compatibilidad hacia atrás. En \texttt{ipc\_protocol.py}, si se recibe un mensaje que no es MessagePack v1, se intenta parsear como JSON plano, permitiendo transiciones graduales.

    \item[Stdout/Stderr Separation] En comunicación inter-proceso, es crítico que:
    \begin{itemize}
        \item Los datos estructurados salgan por \textbf{stdout} en binario (IPC).
        \item Los logs, debug y errores salgan por \textbf{stderr} en texto.
    \end{itemize}
    Esto previene que un \texttt{logging.info()} accidental corrompa el IPC.

    \item[Virtual Threads] Innovación de Java 21 (Project Loom) que permite crear millones de hilos ligeros sin el overhead de hilos del sistema operativo. Ventajas:
    \begin{itemize}
        \item Cada estrategia ejecuta en un hilo virtual separado, sin blocking mutuo.
        \item Código escrito como si fueran threads tradicionales.
        \item Escalabilidad a decenas de estrategias simultáneas.
    \end{itemize}
    Se crean vía \texttt{Executors.newVirtualThreadPerTaskExecutor()}.

    \item[ExecutorService] Pool de hilos en Java que ejecuta tareas de forma asíncrona. En el proyecto:
    \begin{itemize}
        \item \texttt{StrategyRuntimeCoordinator}: ExecutorService basado en Virtual Threads para lanzar y supervisar estrategias de tiempo real.
        \item \texttt{MarketDataService}: ExecutorService para paralelizar descarga de velas.
    \end{itemize}

    \item[Transaccionalidad ACID] Garantía de bases de datos donde:
    \begin{itemize}
        \item \textbf{A (Atomicity):} Transacción se ejecuta completamente o no al todo.
        \item \textbf{C (Consistency):} BD pasa de un estado válido a otro válido.
        \item \textbf{I (Isolation):} Transacciones concurrentes no interfieren.
        \item \textbf{D (Durability):} Cambios persistidos no se pierden tras fallo.
    \end{itemize}
    En el proyecto, MySQL con InnoDB garantiza ACID, permitiendo que AccountingService sea seguro bajo concurrencia.

    \item[Lock Pesimista] Mecánica de base de datos que adquiere un lock exclusivo sobre una fila antes de modificarla, bloqueando otras transacciones. En \texttt{WalletRepository}:
    \begin{verbatim}
    @Lock(LockModeType.PESSIMISTIC_WRITE)
    Optional<Wallet> findByIdWithLock(Long id);
    \end{verbatim}
    Garantiza que ninguna otra transacción lee o modifica la Wallet mientras se procesa.

    \item[Soft Delete (Borrado Lógico)] Patrón donde no se elimina un registro de la BD, sino que se marca como eliminado (columna \texttt{eliminado = 1}). Ventajas:
    \begin{itemize}
        \item Integridad referencial preservada.
        \item Auditoría conserva historial.
        \item Recuperación de datos accidentales es posible.
    \end{itemize}

    \item[Pattern Repository] Patrón de diseño que abstrae el acceso a datos. En el proyecto:
    \begin{itemize}
        \item \texttt{WalletRepository} abstrae consultas a tablas Wallet.
        \item \texttt{VelaRepository} abstrae consultas a tablas Vela.
        \item Spring Data JPA proporciona implementación automática.
    \end{itemize}
    Permite cambiar la base de datos sin tocar la lógica de negocio.

    \item[Facade Pattern] Patrón que proporciona una interfaz simplificada a un subsistema complejo. En el proyecto:
    \begin{itemize}
        \item \texttt{MarketDataService} es una Facade que coordina FetchService e IndicatorsService.
        \item \texttt{PythonBridgeFacade} ordena la orquestación de procesos Python.
    \end{itemize}

    \item[ETL (Extract, Transform, Load)] Proceso de:
    \begin{itemize}
        \item \textbf{Extract:} Obtener datos crudos (API de Binance).
        \item \textbf{Transform:} Limpiar, validar, enriquecer (indicadores técnicos).
        \item \textbf{Load:} Persistir en BD.
    \end{itemize}
    El \texttt{FetchService} implementa ETL incremental optimizado.

    \item[Batch Insert] Inserción masiva de múltiples filas en una sola operación. En MySQL:
    \begin{verbatim}
    INSERT INTO vela (...) VALUES (...), (...), (...)
    ON DUPLICATE KEY UPDATE ...
    \end{verbatim}
    Es 10-100x más rápido que inserts individuales.

    \item[Rate Limiting] Restricción de cuántas peticiones puede hacer un cliente a un servidor por unidad de tiempo. Binance limita a $\approx 1200$ peticiones por minuto. El sistema implementa:
    \begin{itemize}
        \item Esperas proactivas entre peticiones.
        \item Detección reactiva de respuestas 429 (Too Many Requests).
        \item Respeto al header \texttt{Retry-After} cuando está presente.
    \end{itemize}

    \item[Timeout] Límite de tiempo para que una operación se complete. En el proyecto:
    \begin{itemize}
        \item \texttt{STARTUP\_TIMEOUT\_SECONDS = 120:} Máximo para que Python inicie.
        \item \texttt{BATCH\_INACTIVITY\_TIMEOUT\_SECONDS = 300:} Máximo sin actividad.
        \item \texttt{BACKTEST\_INACTIVITY\_TIMEOUT\_SECONDS = 900:} Para backtesting largo.
    \end{itemize}
    Si vence, el proceso se destruye forzosamente.

    \item[Idempotencia] Propiedad de una operación donde ejecutarla múltiples veces produce el mismo resultado que ejecutarla una sola vez. En el proyecto:
    \begin{itemize}
        \item \texttt{ON DUPLICATE KEY UPDATE} en inserciones.
        \item \texttt{correlation\_id} en envelope IPC.
    \end{itemize}
    Permite reintentos seguros sin duplicados.

\end{description}

% ==========================================
% SECCIÓN: TÉRMINOS DE CRIPTOMONEDAS
% ==========================================

\section*{Conceptos de Criptomonedas y Mercados}

\begin{description}
    \item[Symbol (Símbolo de Trading)] Identificador de un par de monedas, ej. BTCUSDT:
    \begin{itemize}
        \item \textbf{BTC:} Bitcoin (activo base).
        \item \textbf{USDT:} Tether en USD (activo cotizante).
    \end{itemize}
    Significa ``precio de 1 Bitcoin en Tether''.

    \item[Timeframe (Marco Temporal)] Período de tiempo que representa cada vela (en formato OHLCV). Ejemplos:
    \begin{itemize}
        \item \textbf{1m:} 1 minuto (para scalping).
        \item \textbf{5m, 15m, 1h:} Comunes en swing trading.
        \item \textbf{1d:} 1 día (para inversión largo plazo).
    \end{itemize}
    El sistema soporta cualquier timeframe que Binance ofrezca.

    \item[Liquidez] Capacidad de comprar o vender una cantidad significativa de un activo sin cambiar drásticamente su precio. Mayor liquidez = más operaciones posibles sin slippage (resbalamiento).

    \item[Slippage (Resbalamiento)] Diferencia entre el precio esperado y el precio real al ejecutar una orden. En mercados ilíquidos, es significativo. El modelo de backtesting del proyecto asume ~0.03\% de slippage fixed.

    \item[Comisión (Fee)] Costo cobrado por el exchange por cada operación. Binance cobra típicamente 0.1\% del volumen. El sistema lo incluye en el cálculo de PnL:
    $$\text{PnL Neto} = \text{PnL Bruto} - \text{Comisiones}$$

\end{description}

% ==========================================
% SECCIÓN: TÉRMINOS DE CRIPTOGRAFÍA Y SEGURIDAD
% ==========================================

\section*{Conceptos de Criptografía y Seguridad}

\begin{description}
    \item[Sal (Salt)] Dato aleatorio único que se añade a una contraseña antes de calcular su hash criptográfico. Cumple tres funciones clave: primero, evita que dos usuarios con la misma contraseña generen hashes idénticos (los atacantes no pueden identificar contraseñas duplicadas); segundo, invalida tablas de arcoíris (rainbow tables) precalculadas, pues cada usuario requeriría una tabla distinta; tercero, en algoritmos como BCrypt se almacena dentro del hash, permitiendo verificar la contraseña sin guardarlo en texto plano. En el sistema, BCrypt genera automáticamente una sal aleatoria por usuario durante el registro, garantizando que no hay dos hashes idénticos aunque las contraseñas sean iguales.

    \item[BCrypt] Algoritmo criptográfico de hash diseñado específicamente para contraseñas \parencite{provos_bcrypt_1999}. A diferencia de funciones de hash rápidas como SHA-256, BCrypt es deliberadamente \textit{lento} y adaptativo: incluye un factor de costo que puede incrementarse sin modificar la base de datos, lo que le permite "envejecer gracefully" mientras la potencia de cálculo aumenta. El factor de costo, junto con la sal aleatoria, hace computacionalmente inviable un ataque de fuerza bruta. En el proyecto, se utiliza en \texttt{UsuarioApplicationService} para cifrar las contraseñas de registro y autenticación.

\end{description}

% ==========================================
% CIERRE DE GLOSARIO
% ==========================================

\vspace{1cm}

\textit{Este glosario es una referencia viva del proyecto. Nuevos términos pueden agregarse conforme se expandan las funcionalidades y se incorporen nuevas técnicas de ML o trading.}